\section{Di\'ametro de un \`arbol y camino}
\whenuse{Hallar el camino m\'as largo de un \`arbol con pesos no negativos.
Con pesos uno, su centro y radio tambi\'en ayudan a conectar varios \`arboles.}
\complexity{Dos recorridos: tiempo y memoria $O(n)$. Son iterativos para
admitir cadenas largas sin agotar la pila.}

\begin{lstlisting}[style=cpp]
pair<ll, vector<int>> tree_diameter(
        const vector<vector<pair<int,ll>>> &tree) {
    int n = tree.size();
    if (n == 0) return {0, {}};
    vector<ll> distance(n);
    vector<int> parent(n);
    auto farthest = [&](int source) {
        fill(distance.begin(), distance.end(), -1);
        fill(parent.begin(), parent.end(), -1);
        vector<int> order{source};
        distance[source] = 0;
        for (int i = 0; i < (int)order.size(); ++i) {
            int v = order[i];
            for (auto [to, weight] : tree[v]) if (distance[to] == -1) {
                assert(weight >= 0);
                distance[to] = distance[v] + weight;
                parent[to] = v; order.push_back(to);
            }
        }
        assert((int)order.size() == n);
        return (int)(max_element(distance.begin(), distance.end()) - distance.begin());
    };
    int start = farthest(0), finish = farthest(start);
    vector<int> path;
    for (int v = finish; v != -1; v = parent[v]) path.push_back(v);
    reverse(path.begin(), path.end());
    return {distance[finish], path};
}
void tree_diameter_example() {
    vector<vector<pair<int,ll>>> tree(4);
    auto add = [&](int a, int b, ll w) {
        tree[a].push_back({b,w}); tree[b].push_back({a,w});
    };
    add(0,1,2); add(1,2,5); add(1,3,4);
    auto [length, path] = tree_diameter(tree);
    cout << length << '\n'; // 9; camino 2--1--3 (o su inverso).
}
\end{lstlisting}

\warning{Requiere un \`arbol conectado, cada arista en ambos sentidos y
distancias que entren en \texttt{long long}. Con pesos negativos, usar DP de
\`arbol en lugar de estos dos recorridos. Para un bosque, procesar componentes.}
Si todas las aristas pesan uno, el radio es $\lceil diametro/2\rceil$ y un
centro es \texttt{path[path.size()/2]}. Es la herramienta base de
Adjoin the Networks (NCPC 2015); no resolver cada componente recorriendo de
nuevo todos los v\'ertices del bosque.
